% Route mod-p: the orthogonal (non-analytic) transcendence route for U via Christol.
% Self-contained section. Banks the maximal-non-automaticity evidence + the verified
% theta-telescoping that explains it, with honest proof status.
\section{The orthogonal route: $(u_n \bmod p)$ and Christol's theorem}
\label{sec:modp}

The transcendence of $U(x)=\sum_{n\ge0}u_nx^n$ over $\mathbb Q(x)$ established in
\S\ref{sec:transcendence} rests on the analytic gate (Lemma~\ref{lem:cos}). We
record here an \emph{independent} line of evidence that touches none of that
machinery, via reduction modulo a prime.

\subsection{Strategy}
Since $u_n\in\mathbb N$ are geodesic-growth counts, $U\in\mathbb Z[[x]]$.
If $U$ were algebraic over $\mathbb Q(x)$, then for all but finitely many primes
$p$ its reduction $U\bmod p$ would be algebraic over $\mathbb F_p(x)$, whence by
\textbf{Christol's theorem} the coefficient sequence $(u_n\bmod p)$ would be
$p$-automatic, i.e.\ its \emph{$p$-kernel}
\[
   \mathcal K_p(u)=\bigl\{\,(u_{p^k n+r})_{n\ge0}\;:\;k\ge0,\ 0\le r<p^k\,\bigr\}
\]
would be \emph{finite}. Contrapositively:
\begin{equation}\label{eq:christol-contra}
   \mathcal K_p(u)\ \text{infinite for one prime }p
   \quad\Longrightarrow\quad
   U\ \text{transcendental over }\mathbb Q(x),\ \text{unconditionally.}
\end{equation}
This needs no analytic input; it is a purely arithmetic/combinatorial criterion.

\subsection{The catalytic block has an explicit $\theta$-telescoping}
The relaxed bulk generating function obeys the catalytic functional
equation~\eqref{eq:catalytic}, which after isolating the dilation $t\mapsto q^2t$
(here $q=x^2$) reads
\begin{equation}\label{eq:modp-firstorder}
   F(q,t)=E(t)+D(t)\,F(q,q^2t),\qquad
   D(t)=\frac{2q(q-1)\,t}{(1-qt)(1-q^2t)},\quad
   E(t)=\frac{2qt}{1-qt}\bigl(1+F(q,q)\bigr).
\end{equation}
Since $F(q,0)=0$, iterating \eqref{eq:modp-firstorder} converges and telescopes.

\begin{proposition}[$\theta$-telescoping of the catalytic block]\label{prop:theta-tel}
With $(a;q)_m=\prod_{i=1}^{m}(1-aq^{i-1})$,
\[
   F(q,t)=\sum_{k\ge0}\alpha_k(t)\,E(q^{2k}t),
   \qquad
   \alpha_k(t)=\prod_{j=0}^{k-1}D(q^{2j}t)
             =\frac{[\,2q(q-1)\,]^{k}\,t^{k}\,q^{\,k(k-1)}}{(qt;q)_{2k}},
\]
where $(qt;q)_{2k}=\prod_{i=1}^{2k}(1-q^{i}t)$. The exponent $q^{k(k-1)}=x^{2k(k-1)}$
is a $\theta$\nobreakdash-/square factor; the Euler product $(qt;q)_{2k}$ in the
denominator (pentagonal-lacunary) redistributes each $\theta$-block over all
degrees.
\end{proposition}

\noindent\emph{Verification.} The product identity for $\alpha_k$ is a finite
$q$-Pochhammer manipulation ($\sum_{j=0}^{k-1}(2j+1)=k^2$ in the numerator power,
$\prod_{j<k}(1-q^{2j+1}t)(1-q^{2j+2}t)=(qt;q)_{2k}$ in the denominator), confirmed
symbolically for $0\le k\le4$; the resulting series reproduces the bulk record
$F(q,1)=2,2,6,2,18,6,42,18,118,50,282,190,706,594$ ($[q^{1}]\dots[q^{14}]$)
exactly. The full true-growth series $U$ is a finite rational assembly (two bulk
blocks, a geometric travel sum, four marker junctions) over $F$, so it inherits
this $\theta$ structure.

\subsection{Computational evidence: maximal non-automaticity}
We computed $(u_n\bmod p)$ by the validated polynomial-time catalytic engine
(a $1$-D collapse of the transfer, $O(N^4)$; reproduces $u_n\bmod p$ exactly for
$n\le42$). The $p$-kernel does not merely grow---it tracks the \emph{full
$p$-ary tree} with $O(1)$ deficit:
\[
\begin{array}{l|llll}
 & k=0 & k=1 & k=2 & k=3\\\hline
 \text{cumulative }|\mathcal K_3|\ (N=180) & 1 & 4 & 13 & 39\\
 \text{full ternary tree} & 1 & 4 & 13 & 40\\[2pt]
 \text{cumulative }|\mathcal K_5|\ (N=130) & 1 & 6 & 30 & -\\
 \text{full quinary tree} & 1 & 6 & 31 & -
\end{array}
\]
(per-level new sections vs.\ maxima: $p=3$, $+1,+3,+9,+26$ vs.\ $1,3,9,27$;
$p=5$, $+1,+5,+24$ vs.\ $1,5,25$). Both counts are robust to the section-overlap
threshold up to the data-length limit ($p=3$ at $k=3$ requires $\ge6$ shared
points per section, which $N=180$ supplies). The subword complexity is likewise
near-maximal ($p=3$: $3,9,27,75$ vs.\ $3^L=3,9,27,81$). This is the strongest
signature finite data affords: the $p$-kernel and factor complexity are as large
as possible, at two independent primes, exactly as
Proposition~\ref{prop:theta-tel} predicts. (By contrast $p=2$ is eventually
periodic---the dilation degenerates---so the route is specific to $p\ge3$.)

\subsection{Honest status}
By Christol, $\mathcal K_p(u)$ infinite $\iff$ $U\bmod p$ transcendental over
$\mathbb F_p(x)$, i.e.\ the explicit partial-$\theta$/Rogers--Ramanujan series of
Proposition~\ref{prop:theta-tel} is non-algebraic mod $p$. Finite kernel data
establishes that $\mathcal K_p(u)$ is \emph{large} (maximal through the deepest
reliable level), not provably \emph{infinite}. The square-support handle is
unavailable: $(u_n\bmod p)$ is dense and base-$p$ pseudorandom (the Euler-product
denominator fills the $\theta$ gaps), so the classical positive-characteristic
gap/zero-set theorems (Kedlaya; Adamczewski--Bell, \emph{Invent.\ Math.}\ 2012)
do not apply directly. Thus this route is genuinely \emph{orthogonal in method}
(automata over $\mathbb F_p(x)$, no analytic gate) and yields unconditional,
maximal evidence, while a proof reduces to the same underlying
$\theta$/square obstruction met by the analytic route and by
Route~D's $x^{k^2}$ specialization.
